%% kapa: exam=1 round=35 year=2024 period=2 form=A subject=appraisal_law subject_ko=감정평가관계법규 seq=19 number=19
%% source: Q-Net 감정평가사 기출문제 (cst00309) artlSeq=5213229 file=raw/1008/5213229/2교시 문제지 원본.pdf
%% conversion: kordoc 4.24.51 -> tools/gen1.py (auto); review=proofread-sonnet; flags=none
\begin{question}{19}
감정평가 및 감정평가사에 관한 법령상 감정평가사에 관한 설명으로 옳지 않은 것은?
\begin{choices}
  \item 감정평가사는 감정평가업을 하기 위하여 1개의 사무소만을 설치할 수 있다.
  \item 견책을 받은 감정평가사는 감정평가사 교육연수의 대상자에 포함된다.
  \item 국유재산을 관리하는 기관에서 5년 이상 감정평가와 관련된 업무에 종사한 사람에 대해서는 감정평가사시험 중 제1차 시험을 면제한다.
  \item 국토교통부장관은 등록한 감정평가사가 파산선고를 받고 복권되지 아니한 경우에는 그 등록을 취소하여야 한다.
  \item 등록한 감정평가사는 5년마다 그 등록을 갱신하여야 한다.
\end{choices}
\end{question}
\answer{2}
